% GENERATED FILE. Edit canonical lessons, then run npm run generate:books.
% canonical-source-hash: fnv1a64:f841b5244b8cf9a4
% canonical-lessons: LA-C394-apparo, LA-C394-dispono, LA-C394-ordino, LA-C394-molior, LA-C394-praevideo

\chapter{Getting Ready}
\label{ch:la-getting-ready}
\hlchaptermodality{394}

\begin{chapteropening}
\textbf{By the end of this chapter:} \emph{I can say that I prepare, set out, put in order and set about a task, and that I foresee what will come.}

Complete the last lesson of chapter 394: I can say that I prepare, set out, put in order and set about a task, and that I foresee what will come.
\end{chapteropening}

\section[apparō, apparāre]{apparō, apparāre — to prepare, to get ready}

\label{lesson:LA-C394-apparo}



\begin{quote}
Before the new one: say the Latin for an undertaking, then the Latin for to try.
\end{quote}

\subsection*{You'll want to know: apparō, apparāre}
\textbf{apparō, apparāre} — \textquotedblleft{}to prepare, to get ready\textquotedblright{}.

\begin{quote}
Convīvium magnum apparāvit.
\end{quote}

\begin{quote}
\emph{He prepared a great banquet.}
\end{quote}

\emph{apparō, apparāre, apparāvī, apparātum}, \emph{ad} + \emph{parō}. \emph{Apparātus}, equipment, splendour; English \emph{apparatus}.

\begin{cousinweb}
\emph{apparō} is built from \emph{ad} \textquotedblleft{}towards\textquotedblright{} and \emph{parō} \textquotedblleft{}to make ready\textquotedblright{} (Lewis \& Short). English \emph{apparatus} comes from its noun.
\end{cousinweb}

\begin{grammarlens}[title={What you've built}]
One more word for telling, explaining and describing at B1.
\end{grammarlens}

\subsection*{Your turn}
\begin{itemize}
\raggedright
  \item \emph{Say it:} \emph{apparō, apparāre}
  \item \emph{Say it:} \emph{apparō, apparāre}, once more
  \item \emph{Say it:} a sentence of your own with \emph{apparō, apparāre}
  \item \emph{Recall:} say the Latin for an undertaking, then the Latin for to try, then say \emph{apparō, apparāre} again
\end{itemize}

\begin{tcolorbox}[breakable,colback=teal!4,colframe=teal!35!black,title={Before you move on}]
What does apparō, apparāre mean? (\textquotedblleft{}To prepare, to get ready\textquotedblright{}.) Say it once more.
\end{tcolorbox}
\section[dispōnō, dispōnere]{dispōnō, dispōnere — to set out, to arrange}

\label{lesson:LA-C394-dispono}



\begin{quote}
Before the new one: say the Latin for to try, then the Latin for to prepare.
\end{quote}

\subsection*{You'll want to know: dispōnō, dispōnere}
\textbf{dispōnō, dispōnere} — \textquotedblleft{}to set out, to arrange\textquotedblright{}.

\begin{quote}
Mīlitēs in mūrīs dispositī sunt.
\end{quote}

\begin{quote}
\emph{The soldiers were posted along the walls.}
\end{quote}

\emph{dispōnō, dispōnere, disposuī, dispositum}: \emph{dis-} + \emph{pōnō}, to place here and there. English \emph{dispose}.

\begin{cousinweb}
\emph{dispōnō} is built from \emph{dis-} \textquotedblleft{}apart\textquotedblright{} and \emph{pōnō} \textquotedblleft{}to put\textquotedblright{} (EtymDB; Lewis \& Short): to put each thing in its own place. English \emph{dispose} comes from it through French.
\end{cousinweb}

\begin{grammarlens}[title={What you've built}]
One more word for telling, explaining and describing at B1.
\end{grammarlens}

\subsection*{Your turn}
\begin{itemize}
\raggedright
  \item \emph{Say it:} \emph{dispōnō, dispōnere}
  \item \emph{Say it:} \emph{dispōnō, dispōnere}, once more
  \item \emph{Say it:} a sentence of your own with \emph{dispōnō, dispōnere}
  \item \emph{Recall:} say the Latin for to try, then the Latin for to prepare, then say \emph{dispōnō, dispōnere} again
\end{itemize}

\begin{tcolorbox}[breakable,colback=teal!4,colframe=teal!35!black,title={Before you move on}]
What does dispōnō, dispōnere mean? (\textquotedblleft{}To set out, to arrange\textquotedblright{}.) Say it once more.
\end{tcolorbox}
\section[ōrdinō, ōrdināre]{ōrdinō, ōrdināre — to set in order, to arrange}

\label{lesson:LA-C394-ordino}



\begin{quote}
Before the new one: say the Latin for to prepare, then the Latin for to set out.
\end{quote}

\subsection*{You'll want to know: ōrdinō, ōrdināre}
\textbf{ōrdinō, ōrdināre} — \textquotedblleft{}to set in order, to arrange\textquotedblright{}.

\begin{quote}
Librōs suōs ōrdināvit.
\end{quote}

\begin{quote}
\emph{He put his books in order.}
\end{quote}

\emph{ōrdinō, ōrdināre, ōrdināvī, ōrdinātum}, built on \emph{ōrdō}, a row. English \emph{ordain} and \emph{coordinate}.

\begin{cousinweb}
\emph{ōrdinō} is built on \emph{ōrdō, ōrdinis} \textquotedblleft{}a row, an order\textquotedblright{} (EtymDB; Lewis \& Short). English \emph{ordain} comes from it through French.
\end{cousinweb}

\begin{grammarlens}[title={What you've built}]
One more word for telling, explaining and describing at B1.
\end{grammarlens}

\subsection*{Your turn}
\begin{itemize}
\raggedright
  \item \emph{Say it:} \emph{ōrdinō, ōrdināre}
  \item \emph{Say it:} \emph{ōrdinō, ōrdināre}, once more
  \item \emph{Say it:} a sentence of your own with \emph{ōrdinō, ōrdināre}
  \item \emph{Recall:} say the Latin for to prepare, then the Latin for to set out, then say \emph{ōrdinō, ōrdināre} again
\end{itemize}

\begin{tcolorbox}[breakable,colback=teal!4,colframe=teal!35!black,title={Before you move on}]
What does ōrdinō, ōrdināre mean? (\textquotedblleft{}To set in order, to arrange\textquotedblright{}.) Say it once more.
\end{tcolorbox}
\section[mōlior, mōlīrī]{mōlior, mōlīrī — to set about (with effort), to undertake; to plot}

\label{lesson:LA-C394-molior}



\begin{quote}
Before the new one: say the Latin for to set out, then the Latin for to set in order.
\end{quote}

\subsection*{You'll want to know: mōlior, mōlīrī}
\textbf{mōlior, mōlīrī} — \textquotedblleft{}to set about (with effort), to undertake; to plot\textquotedblright{}.

\begin{quote}
Magnam rem mōlītur.
\end{quote}

\begin{quote}
\emph{He is setting about a great undertaking.}
\end{quote}

Deponent: \emph{mōlior, mōlīrī, mōlītus sum}, built on \emph{mōlēs}, a mass: to heave a heavy thing. English \emph{demolish} comes from \emph{dēmōlior}, to pull down.

\begin{cousinweb}
\emph{mōlior} is built on \emph{mōlēs} \textquotedblleft{}a mass, a heavy weight\textquotedblright{} (EtymDB; Lewis \& Short): to move a great weight. English \emph{demolish} comes from its compound \emph{dēmōlior}.
\end{cousinweb}

\begin{grammarlens}[title={What you've built}]
One more word for telling, explaining and describing at B1.
\end{grammarlens}

\subsection*{Your turn}
\begin{itemize}
\raggedright
  \item \emph{Say it:} \emph{mōlior, mōlīrī}
  \item \emph{Say it:} \emph{mōlior, mōlīrī}, once more
  \item \emph{Say it:} a sentence of your own with \emph{mōlior, mōlīrī}
  \item \emph{Recall:} say the Latin for to set out, then the Latin for to set in order, then say \emph{mōlior, mōlīrī} again
\end{itemize}

\begin{tcolorbox}[breakable,colback=teal!4,colframe=teal!35!black,title={Before you move on}]
What does mōlior, mōlīrī mean? (\textquotedblleft{}To set about (with effort), to undertake; to plot\textquotedblright{}.) Say it once more.
\end{tcolorbox}
\section[praevideō, praevidēre]{praevideō, praevidēre — to foresee}

\label{lesson:LA-C394-praevideo}



\begin{quote}
Before the new one: say the Latin for to set in order, then the Latin for to set about.
\end{quote}

\subsection*{You'll want to know: praevideō, praevidēre}
\textbf{praevideō, praevidēre} — \textquotedblleft{}to foresee\textquotedblright{}.

\begin{quote}
Perīculum praevīdit et fūgit.
\end{quote}

\begin{quote}
\emph{He foresaw the danger and fled.}
\end{quote}

\emph{praevideō, praevidēre, praevīdī, praevīsum}: \emph{prae-} beforehand + \emph{videō}.

\begin{cousinweb}
\emph{praevideō} is built from \emph{prae-} \textquotedblleft{}before\textquotedblright{} and \emph{videō} \textquotedblleft{}to see\textquotedblright{} (EtymDB; Lewis \& Short).
\end{cousinweb}

\begin{grammarlens}[title={What you've built}]
One more word for telling, explaining and describing at B1.
\end{grammarlens}

\subsection*{Your turn}
\begin{itemize}
\raggedright
  \item \emph{Say it:} \emph{praevideō, praevidēre}
  \item \emph{Say it:} \emph{praevideō, praevidēre}, once more
  \item \emph{Say it:} a sentence of your own with \emph{praevideō, praevidēre}
  \item \emph{Recall:} say the Latin for to set in order, then the Latin for to set about, then say \emph{praevideō, praevidēre} again
\end{itemize}

\begin{tcolorbox}[breakable,colback=teal!4,colframe=teal!35!black,title={Before you move on}]
What does praevideō, praevidēre mean? (\textquotedblleft{}To foresee\textquotedblright{}.) Say it once more.
\end{tcolorbox}
